\documentclass{ximera}
%verbeterd 30/10/26
\title{Serieschakeling}
\author{Daan Lipkens}

\begin{document}
\begin{abstract}
    Eigenschappen van een serieschakeling en de vervangingsweerstand in serie.
\end{abstract}
\maketitle

\section{Serieschakeling}\label{sec:serieschakeling}
\begin{definition}[title={Serieschakeling}]\nl
    In een \textbf{serieschakeling} staan de componenten achter elkaar: er is slechts één pad voor de stroom en er zijn geen vertakkingen.
\end{definition}

\begin{figure}[ht!]
    \captionsetup{width=0.85\textwidth}
    \begin{center}
        \begin{circuitikz}[european]
            \draw (0,3) to[battery1, l=$U$] (0,0) -- (6,0) -- (6,3);
            \draw (0,3) to[R, l=$R_1$] (3,3) to[R, l=$R_2$] (6,3);
        \end{circuitikz}
    \end{center}
    \caption{Twee weerstanden in serie, aangesloten op een spanningsbron.}
    \label{fig:serie}
\end{figure}

\subsection{Stroomsterkte en spanning}
Omdat er geen vertakkingen zijn, kan er geen lading verdwijnen of bijkomen onderweg. Door elke weerstand in figuur~\ref{fig:serie} stroomt dus evenveel lading per seconde. De spanning van de bron wordt verdeeld over de weerstanden.

\begin{proposition}[title={Serieschakeling: stroom en spanning}]\nl
    In een serieschakeling is de stroomsterkte overal gelijk:
    \begin{equation*}
        I = I_1 = I_2 = \dots
    \end{equation*}
    De spanning van de bron is de som van de spanningen over de componenten:
    \begin{equation*}
        U = U_1 + U_2 + \dots
    \end{equation*}
\end{proposition}

\begin{remark}[foldable=true,title={Grootste weerstand, grootste spanning}]\nl
    Door alle weerstanden loopt dezelfde stroom $I$. Uit $U = R \cdot I$ volgt dan dat de weerstand met de grootste waarde ook de grootste spanning krijgt. Een weerstand die twee keer zo groot is, krijgt twee keer zoveel spanning.
\end{remark}

\subsection{Vervangingsweerstand in serie}
We kunnen alle weerstanden in een serieschakeling vervangen door één weerstand, zonder dat de stroom door de bron verandert. Dit noemen we de \textbf{vervangingsweerstand}.

\begin{proposition}[title={Vervangingsweerstand in serie}]\nl
    Voor weerstanden in serie is de vervangingsweerstand $R_{\text{s}}$ de som van de weerstanden:
    \begin{equation*}
        R_{\text{s}} = R_1 + R_2 + \dots
    \end{equation*}
\end{proposition}

\begin{remark}[foldable=true,title={Afleiding voor twee weerstanden}]\nl
    Door beide weerstanden loopt dezelfde stroom $I$, en de bronspanning is de som van beide spanningen:
    \begin{align*}
        U &= U_1 + U_2 \\
        &= R_1 \cdot I + R_2 \cdot I \\
        &= (R_1 + R_2) \cdot I
    \end{align*}
    Een enkele weerstand $R_{\text{s}}$ die dezelfde stroom $I$ geeft bij dezelfde spanning $U$, moet dus voldoen aan $R_{\text{s}} = R_1 + R_2$.
\end{remark}

%---------------------------------------------------------------------------
% VOORBEELD SERIE
\begin{example}[foldable=true,title={Voorbeeld: Twee weerstanden in serie}]\nl
    Een weerstand $R_1 = 10\,\mathrm{\Omega}$ en een weerstand $R_2 = 20\,\mathrm{\Omega}$ staan in serie, aangesloten op een bron van $12\,\mathrm{V}$. Bereken de stroomsterkte en de spanning over elke weerstand.

    \begin{solution}\nl
        \underline{Gegevens:}
        \begin{align*}
            & R_1 = 10\,\mathrm{\Omega} & & R_2 = 20\,\mathrm{\Omega} \\
            & U = 12\,\mathrm{V} & &
        \end{align*}

        \underline{Gevraagd:} $I$, $U_1$ en $U_2$ \\

        \underline{Oplossing:}\\
        \textbf{Stap 1: Vervangingsweerstand}\\
        \begin{align*}
            R_{\text{s}} &= R_1 + R_2 = 10\,\mathrm{\Omega} + 20\,\mathrm{\Omega} = 30\,\mathrm{\Omega}
        \end{align*}
        \textbf{Stap 2: Stroomsterkte}\\
        De stroom door de bron is de stroom door de vervangingsweerstand:
        \begin{align*}
            I &= \frac{U}{R_{\text{s}}} = \frac{12\,\mathrm{V}}{30\,\mathrm{\Omega}} = 0{,}40\,\mathrm{A}
        \end{align*}
        \textbf{Stap 3: Spanning over elke weerstand}\\
        Dezelfde stroom loopt door beide weerstanden:
        \begin{align*}
            U_1 &= R_1 \cdot I = 10\,\mathrm{\Omega} \cdot 0{,}40\,\mathrm{A} = 4{,}0\,\mathrm{V} \\
            U_2 &= R_2 \cdot I = 20\,\mathrm{\Omega} \cdot 0{,}40\,\mathrm{A} = 8{,}0\,\mathrm{V}
        \end{align*}
        Controle: $U_1 + U_2 = 4{,}0\,\mathrm{V} + 8{,}0\,\mathrm{V} = 12\,\mathrm{V} = U$.
    \end{solution}
\end{example}
%---------------------------------------------------------------------------

%---------------
% QUIZ: KERSTVERLICHTING
\begin{question}
    Een oud kerstlichtsnoer bestaat uit lampjes in serie. Wat gebeurt er als één lampje doorbrandt?
    \begin{multipleChoice}
        \choice{Enkel dat lampje gaat uit.}
        \choice[correct]{Alle lampjes gaan uit.}
        \choice{De andere lampjes branden feller.}
        \choice{De andere lampjes branden zwakker, maar blijven aan.}
    \end{multipleChoice}
    \begin{solution}
        Een doorgebrand lampje onderbreekt het enige pad voor de stroom. Er kan nergens nog stroom vloeien, dus alle lampjes gaan uit.
    \end{solution}
\end{question}

\end{document}